a) Der Rand
\mathl{\partial M}{} einer Mannigfaltigkeit mit Rand ist nach
[[Mannigfaltigkeit mit Rand/Rand/Elementare Eigenschaften/Fakt|Fakt   (2)]]
\definitionsverweis {abgeschlossen}{}{,}
das offene Intervall als Teilmenge des $\R^2$ ist aber nicht abgeschlossen.

b) Der Rand
\mathl{\partial M}{} einer Mannigfaltigkeit mit Rand ist nach
[[Mannigfaltigkeit mit Rand/Rand ist Mannigfaltigkeit/Fakt|Fakt]]
eine Mannigfaltigkeit ohne Rand, das abgeschlossene Intervall besitzt aber Randpunkte.

c) Es sei angenommen, dass $\R^2$ eine Mannigfaltigkeit mit dem Rand
\mathl{\R = \R \times \{0\}}{} ist. Zu
\mathl{P \in \R}{} gibt es dann eine offene Umgebung
\mathl{P\in U}{} und eine Homöomorphie
\maabbdisp {\alpha} { U } {V
} {}
mit einer offenen Menge $V \subseteq H$, wobei $H$ eine Halbebene bezeichnet. Dabei können wir $V$ durch eine kleinere offene Halbballumgebung um $P$ ersetzen. Bei einer solchen Karte werden nach
[[Mannigfaltigkeit mit Rand/Rand/Elementare Eigenschaften/Fakt|Fakt   (2)]]
Randpunkte auf Randpunkte abgebildet, d.h. es ist

\mavergleichskettedisp
{\vergleichskette
{\alpha ( U \cap \R  )
}
{ =} { V \cap \delta H
}
{ } {
}
{ } {
}
{ } {
}
}
{}{}{.}
Damit erhält man auch eine Homöomorphie zwischen den Komplementen, also zwischen
\mathkor {} {U \setminus U \cap \R} {und} {V \setminus V \cap \delta H} {.}
Die Halbballumgebung rechts ist aber
\definitionsverweis {wegzusammenhängend}{}{,}
wohingegen die Menge links die disjunkte Vereinigung der Schnitte mit der positiven bzw. der negativen Hälfte ist, die beide offen und auch nicht leer sind, da $U$ eine Ballumgebung von $P$ enthält. Daher ist die Menge links nicht zusammenhängend und es kann keine Homöomorphie geben.

d) wie c).

 [[Kategorie:Latexseite]]
